\documentclass[11pt]{article}
\usepackage[T1]{fontenc}
\usepackage{lmodern,microtype,amsmath,amssymb,booktabs,tabularx}
\usepackage[letterpaper,margin=.86in,top=.8in,bottom=.8in,headheight=14pt]{geometry}
\usepackage{xcolor,fancyhdr,hyperref}
\definecolor{ink}{HTML}{173F47}\hypersetup{colorlinks=true,urlcolor=ink,linkcolor=ink,pdftitle={Two matching theorems for determinants of sums of normal matrices}}
\pagestyle{fancy}\fancyhf{}\fancyhead[L]{\small\color{ink}The Marcus--de Oliveira conjecture}\fancyhead[R]{\small Research continuation}\fancyfoot[R]{\small\thepage}
\setlength{\parindent}{0pt}\setlength{\parskip}{5pt}
\newcommand{\heading}[1]{\section*{\color{ink}#1}}
\newcommand{\C}{\mathbb C}\newcommand{\R}{\mathbb R}\newcommand{\E}{\mathbb E}\newcommand{\pr}{\mathbb P}
\DeclareMathOperator{\co}{conv}\DeclareMathOperator{\rank}{rank}\DeclareMathOperator{\tr}{tr}\DeclareMathOperator{\diag}{diag}
\begin{document}
{\small\sffamily\color{ink}RESEARCH CONTINUATION}\par
{\LARGE\bfseries Two matching theorems for determinants of sums of normal matrices}\par
\heading{1. Statements and scope}
For normal $A,B\in M_n(\C)$ with eigenvalues $a_i,b_j$ counted with multiplicity, set
\[
 z_\pi=\prod_{i=1}^n(a_i+b_{\pi(i)}),\quad
 \mathcal H=\co\{z_\pi:\pi\in S_n\},\quad d=\det(A+B).
\]
MOC asserts $d\in\mathcal H$. All convex coefficients in this note are real and nonnegative.
\par\textbf{Theorem A.} If some $\mu\in\C$ satisfies
\[
 \rank(A-\mu I)+\rank(B+\mu I)\le n+2 ,
\]
then MOC holds for all relative eigenbases. Equivalently an eigenvalue of $A\oplus(-B)$ has multiplicity at least $n-2$ (for $n\ge3$). At equality of the rank sum, the proof constructs a probability law that works simultaneously for all nonzero values of the shifted spectra with fixed eigenvectors. The first rank split beyond the preceding research note is $4+4$ in dimension six.
\par\textbf{Theorem B.} If $n=4$ and one matrix is Hermitian, MOC holds for every normal second matrix. For each fixed Hermitian matrix in the eigenbasis of the second, one probability law works for all complex eigenvalues of the second. By M\"obius covariance the scalar inclusion holds when either spectrum is on a generalized circle, meaning a circle or line. By Shitov's induction, MOC holds for all $n\ge4$ when $A$ is Hermitian and $B$ has at most four nonreal eigenvalues. In particular $n=5$ allows one real eigenvalue of $B$ and four arbitrary complex ones. Theorem B has a finite exact computation in its proof.
\par\textbf{Why these advance different cases.} The rank condition permits both spectra to have nonreal cross-ratios and permits arbitrary eigenbases. The one-Hermitian theorem permits four entirely unrelated eigenvalues on the other side. Shitov [2] established the three-exception result and the general induction theorem. The passage from three to four exceptions here requires the order-four base of Theorem B. Kova\v{c}ec [5] explicitly posed the order-four Hermitian subproblem. The general order-four and order-five problems with both spectra arbitrary remain unresolved here. No claim of historical priority is made.
\par\textbf{Unrestricted order three.} Fix a relative eigenbasis $U$. By the Birkhoff theorem the doubly stochastic matrix $(|U_{ij}|^2)$ is a convex combination of permutation matrices. This law matches all singleton minors. In order three every two-minor is complementary to a singleton, with the same squared modulus by unitarity. The law therefore matches every coefficient in the determinant expansion, proving MOC for arbitrary order-three spectra [1].
\par\textbf{Organization.} Sections 2 and 3 prove Theorem A without computation. Sections 4 and 5 give the projection, finite verification, and analytic inequality for Theorem B. Section 6 classifies spectra for which a single law works for every partner. Sections 7 and 8 treat open sets, matrix-algebra bundles, and a family with unrestricted spectra. Section 9 gives an exact obstruction at one fixed pair of spectra. Verification instructions are in Section 10.
\newpage
\heading{2. The two-frame matching theorem}
Let $E\in\C^{r\times2}$ and $F\in\C^{s\times2}$ have orthonormal columns. Write $e_i,f_j$ for their rows, and $[v,w]=v_1w_2-v_2w_1$. There is a random matching of two edges, with distinct endpoints in each set, for which
\[
 \begin{aligned}
 \pr(\text{row set}=\{i,h\})&=|[e_i,e_h]|^2,\\
 \pr(\text{column set}=\{j,l\})&=|[f_j,f_l]|^2,\\
 \pr((i,j)\text{ is an edge})&=|[e_i,f_j]|^2 .
 \end{aligned}                                                   \tag{1}
\]
Both frames can have arbitrarily many rows. Their selected pairs may depend on one another.
\par\textbf{Convert determinants into squared distances.} To each nonzero row $v$ associate its mass $p=\|v\|^2$ and its unit Bloch vector $n\in\R^3$ by $v^*v=\frac p2(I+n\cdot\sigma)$, with $\sigma$ the three Pauli matrices. A zero row can be omitted. In either frame $\sum p=2$ and $\sum pn=0$. For any two rows from either frame, including one from each,
\[
 |[v,w]|^2=\|v\|^2\|w\|^2-|vw^*|^2=\tfrac14 p_v p_w\|n_v-n_w\|^2. \tag{2}
\]
\textbf{Choose axes with equal variances.} Form the difference $H$ of the two real covariance sums $\sum p nn^T$. It is symmetric and traceless. There is an orthonormal basis $t_1,t_2,t_3$ with $t_l^THt_l=0$ for all $l$. Indeed a traceless real symmetric form has a zero on the unit sphere, and its restriction to the perpendicular space again has trace zero. Denote the common second moment on axis $t_l$ by $V_l$. Then $\sum_l V_l=2$.
\par On each axis with positive $V_l$, replace the row $v$ in either frame by the real row
\[
 \widetilde v_l=\left(\sqrt{p_v/2},\ \sqrt{p_v/V_l}\,(n_v\cdot t_l)\right).
\]
The two columns remain orthonormal, by the mass, mean, and second moment identities. The squared wedge of two new rows is $p_vp_w((n_v-n_w)\cdot t_l)^2/(2V_l)$. Equation (2) therefore expresses the entire list of squared determinants as a convex combination of the three new lists, with weights $V_l/2$. A zero variance contributes no distance on that axis.
\par\textbf{Determine the extreme points.} The set of lists from all frame pairs of fixed sizes is compact, and so is its finite-dimensional convex hull. Each extreme point is realized by frames. Its convex decomposition above forces each positive-weight projected list to equal that extreme point. The list determines each row mass since $\sum_h |[e_i,e_h]|^2=\|e_i\|^2$, and likewise for $F$. Equality of the masses before and after projection gives $(n_v\cdot t_l)^2=V_l/2$ on all active rows. Square roots are taken consistently in both frames. Every projected row is now parallel to $(1,1)$ or $(1,-1)$. Each direction has total squared norm one in each frame.
\par Independently choose one label from each direction of each frame with probabilities given by these squared norms. Join labels on opposite directions across frames. Wedges of parallel rows vanish, and the squared wedge of orthogonal rows equals the product of their squared norms. This proves (1) at each extreme point. Taking mixtures proves (1) for all frames. Conversely each individual matching is itself a frame list, by taking two unit coordinate directions and setting unused rows to zero.
\newpage
\heading{3. Apply the matching to the original determinant}
Shift the matrices by opposite scalars to get ranks $r,s$. Write their nonzero eigenvalues as $u_i,v_j$, and use isometries $X,Y$ of corresponding eigenvectors. Write
\[
 T=[X\mid Y],\quad D=\prod_i u_i\prod_j v_j,\quad x_i=1/u_i,\quad y_j=1/v_j,
 \qquad \Lambda=\diag(u_1,\ldots,u_r,v_1,\ldots,v_s).
\]
Then $A+B=T\Lambda T^*$. If $T$ lacks full row rank, $d=0$. Both supports are proper and a zero-zero pairing gives a zero vertex. Assume full row rank and put $k=r+s-n$. The common support has dimension $k$. Let the columns of $W$ form an orthonormal basis of it. Then $E=X^*W$, $F=Y^*W$ are isometries. The columns of $2^{-1/2}\binom E{-F}$ form an orthonormal basis of $\ker T$. Cauchy--Binet on $(TT^*)^{-1/2}T$ and complementary minors of its unitary completion give
\[
 \det(A+B)=Dq\det(E^*D_xE+F^*D_yF),\qquad q=2^{-k}\det(TT^*)\le1. \tag{3}
\]
In fact $TT^*=XX^*+YY^*$ equals $2I$ on the intersection. Its positive restriction to the complement has trace equal to the dimension, so its determinant is at most one. If either support is full, $q=1$.
\par For $k=2$, Cauchy--Binet on the stacked frames expands the small determinant in (3) into pure $x_ix_h$, pure $y_jy_l$, and mixed $x_iy_j$ terms. Their respective coefficients are exactly the three probabilities in (1). If the matching is $(i,j),(h,l)$, take the product
\[
 (x_i+y_l)(x_h+y_j).
\]
The two mixed terms correspond to the actual edges. The expectation is the small determinant. After multiplying by $D$, this product is an original permutation vertex: pair $u_i$ with $v_l$ and $u_h$ with $v_j$. Assign all remaining nonzero entries to zeros on the opposite side. Multiply the law by $q$. If $q<1$, a zero vertex absorbs the remaining weight.
\par For $k=1$ the small determinant is the sum of two convex averages, one in $x_i$ and one in $y_j$. Sample their labels independently, and use the vertex $D(x_i+y_j)$. The factor $q$ is treated as above. For $k=0$ the small determinant is 1, and the vertex $D$ suffices with the same possible zero weight. A negative rank excess makes every sum singular. This completes Theorem A.
\par\textbf{Exact example.} The file \texttt{verify\_rank\_sum\_six.py} constructs two normal $6\times6$ matrices of rank four. All four row directions are distinct in each compressed frame. The verifier checks 28 coefficients using 20 matchings, applies $q=1296/4225$, and assigns the remaining mass to a zero vertex. Direct evaluation and the resulting 21 permutation weights both give
\[
 \det(A+B)=\frac{-225594+186786i}{4225}.
\]
\par\textbf{Interpretation on a qubit.} Equation (1) gives a local joint probability model for a two-qubit singlet when both parties may choose between the same two arbitrary rank-one measurements. A random order of the matching edges $(i,j),(h,l)$ supplies respective pairs of preassigned outcomes $(i,l)$ and $(h,j)$. The Born pair probabilities are the squared wedges divided by two. Three shared measurements can already fail. For projective axes with pairwise dot product $-1/2$, equal measurements force opposite signs across parties. A local model would give signs $X,Y,Z$ with $\E(XY+YZ+ZX)=-3/2$, whereas $XY+YZ+ZX\ge-1$ for any three signs.
\newpage
\heading{4. Four eigenvalues on one line}
Place the arbitrary spectrum on the rows and a real spectrum $b$ on the columns. For a relative unitary $U$, let $m_{I,J}=|\det U[I,J]|^2$ and let $P_{I,J}(\pi)$ be the indicator that $\pi$ maps $J$ onto $I$. The determinant expansion is
\[
 d=\sum_{|I|=|J|} a_{I^c} b_J m_{I,J},\qquad b_J=\prod_{j\in J}b_j .
\]
A probability law matching every singleton also matches every triple by complementation. For a fixed two-set $I$, let $d_J$ be the difference between the actual $m_{I,J}$ and its proposed expectation. Cauchy--Binet gives $\sum_{J\ni j}m_{I,J}=\sum_{i\in I}m_{\{i\},\{j\}}$, with the identical relation for each permutation. For an explicit check, write the two rows as $u,v$. Sum $|u_jv_k-u_kv_j|^2$ over $k$ and use row orthogonality. Thus $\sum_{J\ni j}d_J=0$. On four labels these equations say $d_J=d_{J^c}$ and the three disjoint-complement classes have values summing to zero.
\par It follows that $\sum b_Jd_J$ depends only on the three real coefficients $b_i b_j+b_k b_l$ for the three partitions $ij\mid kl$. Subtract their minimum and divide by the sum of the remaining two. Relabel columns. It is sufficient to match, for all $I$,
\[
 p\, m_{I,\{4,2\}}+(1-p)m_{I,\{4,3\}},\qquad 0\le p\le\tfrac12, \tag{4}
\]
along with the 16 singleton entries. Equal coefficients require only the singleton law. This reduction also explains why the second spectrum can be arbitrary complex: the law works before evaluation of any $a_i$.
\par\textbf{The finite projected cone.} There are 24 permutation profiles. In their full graded-minor coordinates the augmented span has dimension 14. Its cone of nonnegative affine costs has 82 extreme rays: 16 singleton coordinates, 18 two-minor coordinates up to complementation, and 48 row/column relabelings of the mixed cost
\[
 m_{3,3}+m_{2,4}+m_{12,13}-m_{13,23}.
\]
Here $m_{12,13}$ uses two-element index sets. The supplied integral boundary certificate is included and rechecked: 9050 simplices lie in these supporting facets. Their oriented boundary cancels, and all 6846 nonzero simplex volumes have the same sign relative to an interior point. The resulting nonzero closed top chain covers the connected boundary sphere, certifying completeness.
\par The span of the singleton costs and (4) has dimension 12 for generic $p$. It has codimension two in the full space. Restrict the 82-ray cone to that subspace. A minimal nonnegative dependence of the two quotient coordinates has support at most three. A positive sum of rays vanishes only on common zero vertices. To give a facet in this projected space it must vanish on at least 11 permutations. The exact enumeration checks all one-, two-, and three-ray supports meeting this contact bound. Supports that use only ordinary minors already give nonnegative costs on $U$. After relabeling, every surviving mixed case follows from the inequality on the next page. The enumeration uses polynomial signs in $r=p/(1-2p)$, not a sample of $p$ values. Generic validity suffices at exceptional parameters and endpoints by compactness of the probability simplex and continuity of (4).
\newpage
\heading{5. The inequality behind the 96 mixed cases}
All linear identities between permutation coordinates hold on the unitary minor profile. Indeed the real signed weights $c_\pi=\Re(\operatorname{sgn}(\pi)\prod_j U_{\pi(j),j}/\det U)$ have sum one. Grouping the determinant expansion over $\pi(J)=I$ and applying the complementary-minor formula gives $\sum_{\pi(J)=I}c_\pi=m_{I,J}$. This transfers the finite calculation to unitary profiles.
\par Write $Q_0=m_{34,34}$, $Q_1=m_{23,23}$, $V=m_{24,13}$, $L=m_{2,3}+m_{3,2}$, and $F=L+Q_0-V$. The mixed costs, up to relabeling, are $w_0Q_0+w_1Q_1+w_f F$, with $w_0,w_f>0$, $w_1\ge w_f$ and $w_0(w_1-w_f)=w_f^2$. It suffices to prove for real orthogonal matrices
\[
 V\le L+(\sqrt{Q_0}+\sqrt{Q_1})^2. \tag{5}
\]
Indeed each cost then exceeds or equals $(\sqrt{w_0Q_0}-\sqrt{(w_1-w_f)Q_1})^2$.
\par\textbf{Reduction to real orthogonal matrices.} Multiply columns by phases to make one anchor row a real unit vector $v$. On the real space $v^\perp$ the other rows give real symmetric positive effects $H_i=\Re(u_i^*u_i)$, each of trace one and with sum the identity. Any extreme triple has a rank-one member. Otherwise symmetric perturbations supported on the three ranges have at least nine dimensions in total. The sum imposes at most six conditions and the three traces add at most two more. A nonzero perturbation in both signs contradicts extremality. The rank-one member is an orthogonal projection. The two others sum to the identity on its two-dimensional complement, commute, and have trace one. Their simultaneous diagonalization explicitly gives a mixture of pairs of complementary rank-one projections. Each singleton coordinate and each square of an anchored minor is affine in the effects, by expansion against the real row $v$. All other pair coordinates follow by complementary minors. Thus every complex profile is an average of real orthogonal profiles. It suffices to average the resulting linear cost inequalities.
\par\textbf{Quaternion proof of (5).} Row signs permit determinant one. Represent the rotation by left and right multiplication by unit quaternions $q,r$ in the basis $1,i,j,k$. Put $P_i=q_i^2$, $R_i=r_i^2$, $s_i=P_i+R_i$ ($0\le i\le3$). The signed minors for $Q_0,Q_1,V$ are, respectively,
\[
 d_0=s_0+s_1-1,\qquad d_1=1-s_1-s_2,\qquad v=R_0+R_2-P_0-P_2.
\]
Thus $(|d_0|+|d_1|)^2=\max((s_0-s_2)^2,(s_1-s_3)^2)$. Direct multiplication gives $L=2[(q_0r_3-q_3r_0)^2+(q_1r_2+q_2r_1)^2]\ge2(D_0^2+D_1^2)$, where $D_0=\sqrt{P_3R_0}-\sqrt{P_0R_3}$ and $D_1=\sqrt{P_1R_2}-\sqrt{P_2R_1}$. Interchanging $P,R$ preserves the bounds and allows $v\ge0$. Put $u=R_2-P_0$, $h=R_0-P_2$, $w=P_3-R_1$, $z=P_1-R_3$. Then $u+h=w+z=v$. The differences of $v^2$ from the two squared terms in the maximum are $4uh$ and $4wz$. If either is nonpositive the result follows. Otherwise $u,h,w,z>0$. Apply $\sqrt{(a+b)(c+d)}\ge\sqrt{ac}+\sqrt{bd}$ to $(R_1+w)(P_2+h)$ and $(R_3+z)(P_0+u)$. This gives $D_0+D_1\ge\sqrt{wh}+\sqrt{zu}\ge2(uhwz)^{1/4}$. Hence $2(D_0^2+D_1^2)\ge4\min(uh,wz)$, as required.
\newpage
\heading{6. Which order-four spectra admit a law for every partner?}
\textbf{Theorem C.} Fix a list $b$ of four complex numbers. The following are equivalent.
\begin{enumerate}
\item For each $U\in U(4)$ there is a probability law on $S_4$ giving the MOC identity for $D_a+U D_b U^*$ simultaneously for every $a\in\C^4$.
\item All four points of $b$ lie on a generalized circle.
\end{enumerate}
This is a sharp statement about the order of the quantifiers. Ordinary MOC permits a different law for every pair of spectra.
\par\textbf{Sufficiency and the higher-dimensional consequence.} In the difference convention set $A_0=D_a$, $C=-U D_b U^*$ and $g(z)=(\alpha z+\beta)/(\gamma z+\delta)$. Multiplication gives
\[
 (\gamma A_0+\delta I)(g(A_0)-g(C))(\gamma C+\delta I)=(\alpha\delta-\beta\gamma)(A_0-C).
\]
The determinant and every permutation product therefore acquire the same multiplier. Choose $g$ taking the circle of $-b$ to the real line, with a pole away from these points. The transformed eigenvectors of $C$ stay fixed. By Theorem B its weights work for all choices of $g(a)$ and hence for all $a$ away from the pole. Polynomial continuity includes every remaining $a$. Repeated eigenvalues give at most three distinct points, which always lie on a generalized circle. For $n\ge4$, Shitov's Theorem 19 [2] lifts the order-four Hermitian case to all pairs with one Hermitian matrix and at most four nonreal eigenvalues of its partner. The same transformation gives a sufficient condition that one generalized circle contain every $a_i$ and at least $n-4$ of the reflected points $-b_j$, counting multiplicity.
\par\textbf{Necessity.} Suppose the four points are distinct and nonconcyclic, and that one law works for a fixed $U$ and every $a$. By varying one diagonal variable and setting the other three equal, the polynomial identity recovers the characteristic polynomial of each principal three-dimensional deletion of $U D_b U^*$. The cofactor expansion gives, for each $i$,
\[
 \sum_j\big(m_{i,j}-\pr[\pi(j)=i]\big)\prod_{l\ne j}(t-b_l)=0.
\]
Evaluation at $t=b_j$ gives the singleton marginals. For fixed output pair $I$, equality of the coefficient of $a_{I^c}$ gives $\sum_J b_Jd_J=0$ in the notation of Section 4. The residuals on complementary pairs coincide and their three values $u,v,w$ sum to zero. The extra equation is $s_1u+s_2v+s_3w=0$, where $s_1=b_1b_2+b_3b_4$, $s_2=b_1b_3+b_2b_4$, $s_3=b_1b_4+b_2b_3$. The ratio $(s_1-s_2)/(s_1-s_3)$ is a cross-ratio of $b$ and is nonreal. These three complex points are affinely independent over $\R$. Hence $u=v=w=0$. The law must reproduce every squared minor of $U$.
\par Such a law fails for some $U$ [3]. There is also an explicit check in the supplement: with indices from zero, $m_{0,1}+m_{1,3}-m_{02,12}+m_{03,23}$ is nonnegative at all permutations and equals $-2048/203125$ at the rational orthogonal matrix in \texttt{local\_pole\_theorem.md}. This proves necessity and identifies the exact obstruction to a universal-in-$a$ argument.
\newpage
\heading{7. Open sets of spectra and a $C^*$-algebra formulation}
\textbf{An explicit neighborhood valid for every relative eigenbasis.} Consider
\[
 a=(-3,-1,2,4),\qquad b=(6-5i,-6+3i,2+4i,-7).
\]
Each of the eight entries may move independently by a complex displacement of modulus at most $10^{-8}$. Every pair in this closed polydisc satisfies MOC for every $U\in U(4)$. In particular changing $a_1$ alone to $-3+10^{-9}i$ gives two simple spectra with nonreal cross-ratios. All eight points of the combined list $a,-b$ are distinct.
\par Here are the estimates behind the certification. For an oriented polygon edge $z_pz_q$ write $v_j=\Im(\overline{z_q-z_p}(z_j-z_p))$. Exact integer arithmetic verifies that the polygon has twelve edges, all joining permutations differing by a single transposition, with every nonendpoint slack at least $\delta=204$. Fix one edge. Let $S$ be the union of the endpoint permutation supports and $M=\sum_{(i,j)\notin S}|U_{ij}|^2$. At the center the nonnegative law of Section 4 has singleton marginals; its mass outside the two endpoints is at least $M/4$. Thus the determinant edge slack is at least $\delta M/4$. In the real signed law of Section 5, each excluded permutation uses at least two entries outside $S$, so $|c_\pi|\le M/2$. If the edge slacks change by at most $e$, the new determinant slack is at least $M(\delta/4-11e)$, uniformly over all $U$. The squared vertex diameter is $25,457,650$. Telescoping the four factors bounds each vertex displacement by $R=18000\cdot10^{-8}$. The oriented area changes by at most $4R(6000+R)<204/44$. All twelve edge inequalities persist. The supplement gives the exact check and a stability theorem for simple-spectrum centers even when their edges have other permutation types.
\par\textbf{Continuous central weights.} Let $X$ be compact Hausdorff and let $\mathcal A$ be the algebra of continuous sections of a locally trivial bundle of $M_4(\C)$ algebras over $X$. Suppose $P_1,\ldots,P_4$ and $Q_1,\ldots,Q_4$ are two continuous ordered resolutions of the identity by rank-one projections. Fix four scalars $b_j$ on one generalized circle. There exist nonnegative central functions $\lambda_\pi\in C(X)$ summing to one such that, for every choice of $a_i\in C(X)$, the identity of central functions
\[
 \det\!\left(\sum_i a_iP_i+\sum_j b_j Q_j\right)
   =\sum_{\pi\in S_4}\lambda_\pi\prod_i(a_i+b_{\pi(i)}) \tag{6}
\]
holds. The determinant is the ordinary determinant in each fiber. Local frames may be used to compute it.
\par To prove continuity, first take $b$ real. The projected coordinates (4) and the singleton coordinates form a continuous map into the fixed permutation polytope of Section 4. Indeed the squared minors are traces of products of the rank-one projections onto the relevant exterior-power lines. This description is invariant under changes of local bases and phases. Triangulate the polytope using its vertices, for example by pulling triangulation. Barycentric coordinates on its simplices, extended by zero at unused vertices, agree on shared faces. They define continuous weights on the polytope with at most twelve nonzero entries at each point, since its affine dimension is at most eleven. Composing them with the profile gives (6). The common multiplier from Section 6 transports the same central weights to concyclic $b$. Values at a pole follow by the polynomial identity in each fiber. No global eigenvector choice is required. The existence of the two ordered resolutions is an explicit hypothesis. The scalar order-four result also applies to $sA+tB$ for all complex weights $s,t$ whenever one original spectrum lies on a generalized circle, by applying Theorem B to the scaled matrices.
\newpage
\heading{8. An algebraic-geometric route with unrestricted spectra}
\textbf{Theorem D.} In order four the ordinary scalar MOC holds for arbitrary complex $a,b$ whenever the relative eigenbasis is $U=\operatorname{Sym}^3(g)$ for $g\in SU(2)$, in the orthonormal binary cubic basis, up to row and column phases and permutations. The proof and exact certificate appear in the \texttt{spin} supplement.
\par\textbf{Why this case tests a different argument.} Write $t=|g_{12}|^2$. At $t=2/5$ a valid inequality of the common-minor permutation polytope takes the strictly negative value $-7/125$. Hence a single law reproducing all minors cannot establish this family. The scalar conclusion still holds for every pair $a,b$.
\par\textbf{Mapping to eight points.} Regard the two spectral lists as eight points in the projective line. The products $z_\pi$ are products of brackets joining the two lists. Their invariant linear span has dimension fourteen. A putative scalar violation would give a real linear support cost with nonnegative values on every permutation and a negative value at the unitary profile. Such costs are real parts of complex product coordinates, with an arbitrary real constant added. Before the constant is added they lie on the secant variety of the eight-point quotient. Howard, Millson, Snowden and Vakil [4] prove that the Hessian of its skew cubic is a square modulo that cubic. Eliminating the constant from this relation gives a homogeneous degree-nine polynomial $P$ that vanishes on every realizable cost. The supplement proves the elimination and normalization, including costs where the intermediate denominator vanishes.
\par On the interval where the common law fails, a violating cost would have to lie in a particular 13-dimensional simplex, with positive barycentric coordinate at one designated vertex. The exact calculation expands $P$ in its degree-nine Bernstein basis. Of its 497,420 coefficients, 480,297 are positive and 17,123 vanish. The coefficient at the designated vertex is positive. Each basis function is nonnegative on the simplex, and the designated one is strictly positive throughout the part in question. Thus $P$ cannot vanish at a violating cost. This is an algebraic exclusion of every pair of spectra for this family. The remaining spin parameters have an ordinary common-minor law.
\par\textbf{An exact target for unrestricted order four.} After a common M\"obius transformation the first three entries of the first spectrum can be made real. Fix them and fix the entire second spectrum and $U$. Leave the fourth entry $z$ free. Write $d=d_0+zd_1$ and $z_\pi=v_{0\pi}+zv_{1\pi}$ by multilinearity. Then MOC for every $z\in\C$ is equivalent to the four-real-dimensional inclusion
\[
 (d_0,d_1)\in\co\{(v_{0\pi},v_{1\pi}):\pi\in S_4\}\subset\C^2. \tag{7}
\]
Indeed any real separating functional has form $\Re(\eta x+\xi y)$. If $\eta\ne0$, set $z=\xi/\eta$; if $\eta=0$, a strict separating functional persists for nearby nonzero $\eta$. Conversely any convex representation in (7) works for all $z$. For $z\in\R$ the order-four theorem already proves all the corresponding support inequalities. Complex $z$ introduces the missing directions. Neither the common-law method of Section 6 nor the spin certificate currently covers all of (7).
\newpage
\heading{9. Even at fixed spectra the singleton constraint can fail}
Equation (7) offers a smaller convex test. One might try to retain the singleton marginals $\pr[\pi(i)=j]=|U_{ij}|^2$ used in the Hermitian proof. The following example rules out this requirement even for a single pair of fixed, simple spectra. In this section $\pi$ maps a row label to a column label, and the labels start at zero.
\[
 U=\frac1{14}\begin{pmatrix}-7&7&7&7\\-11&-7&-5&1\\1&7&-11&5\\-5&7&-1&-11\end{pmatrix},\qquad
 a=\left(0,1,2,\frac{8+2i}{17}\right),\qquad b=(1,2,3,i).
\]
The rows of $U$ are orthonormal. Neither quadruple lies on a generalized circle, since each contains three distinct real points and one nonreal point. No $a_i$ equals a reflected $-b_j$. Direct evaluation gives
\[
 d=\frac{7233}{833}+\frac{70745}{3332}i
   =\frac{16}{245}z_{0321}+\frac{257}{588}z_{1302}+\frac{209}{420}z_{3120}. \tag{8}
\]
These positive weights sum to one. To exclude every scalar law that also matches the singleton probabilities, take $\eta=3-5i$ and the real array
\[
 C=\begin{pmatrix}-108&-128&-144&-108\\-24&-9&-4&0\\-32&0&12&0\\0&0&0&0\end{pmatrix}.
\]
At all 24 permutations the cost $\Re(\eta z_\pi)+\sum_i C_{i,\pi(i)}$ belongs to $\{0,4,8\}$. But
\[
 \Re(\eta d)+\sum_{i,j}C_{ij}|U_{ij}|^2=-6/49.
\]
An average of the 24 nonnegative costs cannot give this negative value. The verifier checks the separation and (8) exactly.
\par\textbf{Why such a witness can be found.} Write $a=(x,y,w,z)$ with $x,y,w$ distinct real and $z$ free. Choose a nonconcyclic $b$. If for every $z$ there were a scalar law with the correct singleton marginals, separation as in (7) would give one law with those marginals working for all $z$. Here the fixed simplex is replaced by its compact convex intersection with the singleton constraints.
\par For that law the paired-minor residual $d_{I,J}$ has equal values on complementary row pairs and also on complementary column pairs, by the incidence identities of Section 4 on both sides. Write the resulting real $3\times3$ residual matrix as $D$; all its row and column sums vanish. With $s(b)$ the three scores from Section 6, matching the determinant requires $s(a)^TDs(b)=0$. The constant and linear coefficient vectors of $s(a)$ in $z$ are $(xy,xw,yw)$ and $(w,y,x)$. Together with $(1,1,1)$ they are linearly independent, by the Vandermonde determinant for $x,y,w$. Thus $Ds(b)=0$. Since $b$ is nonconcyclic, the three real vectors $\mathbf1,\Re s(b),\Im s(b)$ span $\R^3$. Therefore $D=0$, so this law would reproduce the entire minor profile. At any mixer failing the common-minor test, some fixed $z$ must violate the scalar law with prescribed singleton probabilities. Equation (8) shows exactly how the ordinary scalar inclusion can survive.
\newpage
\heading{10. Verification and sources}
\textbf{How to reproduce the computations.} All supplied certificates use integer or rational arithmetic. The rank theorem in Sections 2 and 3 has an independent proof on paper. Run the following from \texttt{Code/} in the repository. Python checks must run with assertions enabled.
\begin{itemize}
\item \texttt{python3 moc4\_hermitian\_certificate.py} verifies the 82-ray face lattice and all 96 surviving mixed circuits as polynomial identities. Python 3.10 or newer suffices.
\item \texttt{python3 verify\_cone82.py} checks the supporting facets using the independently supplied boundary chain in \texttt{../data/cone82.json}. A third check compiles \texttt{audit\_82\_exact.cpp} with a C++17 compiler and runs it with assertions enabled. This checks the signs of 1,961,256 determinants and scans 2,496,144 possible thirteen-vertex supports.
\item For the rank-four example run \texttt{python3 verify\_rank\_sum\_six.py}. For the rational obstruction in Section 6 run \texttt{python3 verify\_local\_pole.py}. Both require SymPy.
\item \texttt{python3 verify\_moc4\_open\_spectra.py} checks every numerical assertion in the open-set example. This uses the Python standard library.
\item \texttt{python3 verify\_singleton\_obstruction.py} checks the fixed-spectrum separation and the three positive weights in (8), using SymPy.
\item In the \texttt{spin} directory, \texttt{python3 verify\_scalar.py} regenerates the invariant and its exact Bernstein coefficients. It requires SymPy, a C++17 compiler with OpenMP, and GMP. \texttt{python3 check\_coverage.py} checks an additional rational orthogonal example.
\end{itemize}
\par\textbf{Scope of review.} The finite checks have executable sources. The arguments linking the inequalities to the conjecture are mathematical proofs presented for independent review. Unrestricted order three is already known. The first general unresolved case in this work has order four and two spectra off generalized circles. The all-dimensional rank theorem includes, in order four, any repeated eigenvalue of $A\oplus(-B)$, and in order five, any eigenvalue of $A\oplus(-B)$ of multiplicity at least three.
\par\textbf{References.}
\par [1] N. Bebiano and J. P. da Provid\^encia, ``Revisiting the Marcus--de Oliveira Conjecture,'' \emph{Mathematics} 13 (2025), 711. \href{https://doi.org/10.3390/math13050711}{doi:10.3390/math13050711}.
\par [2] Y. N. Shitov, ``A further sufficient condition for the determinantal conjecture,'' \emph{Izvestiya: Mathematics} 89 (2025), 862--869. Theorem 1 gives the three-exception result. Theorem 19 gives the induction used here. \href{https://doi.org/10.4213/im9292e}{doi:10.4213/im9292e}.
\par [3] S. W. Drury, ``The determinantal conjecture and Hadamard type inequalities,'' \emph{Linear Algebra and its Applications} 246 (1996), 279--297. The four-dimensional common-law obstruction concerns the weights on minors. \href{https://doi.org/10.1016/0024-3795(94)00361-0}{Publisher record}.
\par [4] B. Howard, J. Millson, A. Snowden, R. Vakil, ``The geometry of eight points in projective space: representation theory, Lie theory and dualities,'' Proposition 3.19 and the paragraph following it. \href{https://arxiv.org/abs/1103.5255}{arXiv:1103.5255}.
\par [5] A. Kova\v{c}ec, ``The Marcus--de Oliveira conjecture, bilinear forms, and cones,'' \emph{Linear Algebra and its Applications} 289 (1999), 243--259. \href{https://doi.org/10.1016/S0024-3795(97)10004-0}{Publisher record}.
\end{document}
